% PIC-tutor native LaTeX book — listings, line wrapping, source-path notation.
%
% pygments is not installed on the reference machine (Phase 0 record), so
% code is set with fvextra Verbatim (breaklines) instead of minted. Raw code
% is written verbatim; underscores and backslashes inside code stay literal.

% Inline code: `code` in the Markdown source.
% \code: token may contain spaces; rendered literally (no line breaking).
% \cpath: token has no spaces; rendered via url's \path so long identifiers
%         (e.g. WarpX::PushParticlesandDeposit(), test_3d_electrostatic_...)
%         may break at punctuation instead of overflowing the line.
\newcommand{\code}[1]{\texttt{\detokenize{#1}}}
\newcommand{\cpath}[1]{\path{#1}}

% Fenced code blocks. breakanywhere is risky for CJK but acceptable for the
% mostly-ASCII code in this book; breaklines alone handles long lines.
\DefineVerbatimEnvironment{codeblock}{Verbatim}{%
  breaklines,breakanywhere,
  fontsize=\small,
  frame=single,
  rulecolor=\color{PicTutorCodeBorder},
  framesep=4pt,
  baselinestretch=0.95,
}

% Bash console blocks: same body, visually marked as shell input/output.
\DefineVerbatimEnvironment{consoleblock}{Verbatim}{%
  breaklines,breakanywhere,
  fontsize=\small,
  frame=single,
  rulecolor=\color{PicTutorCodeBorder},
  framesep=4pt,
  baselinestretch=0.95,
}
